\documentclass[10pt, a4paper]{article}
\usepackage[utf8]{inputenc}
\usepackage[T1]{fontenc}
\usepackage[italian]{babel}
\usepackage{geometry}
\geometry{a4paper, margin=2.5cm}
\linespread{1.2}
\usepackage{svg}
\usepackage{float}
\usepackage{graphicx}
\usepackage{amsmath}
\usepackage{hyperref}
\usepackage{xcolor}
\usepackage{listings}
\usepackage{tcolorbox}
\tcbuselibrary{listings}
\usepackage{sourcecodepro}
\usepackage{booktabs}
\usepackage[normalem]{ulem}

\usepackage{fancyhdr}
\setlength{\headheight}{15pt} 
\pagestyle{fancy}
\fancyhf{}
\fancyhead[R]{\footnotesize\color{gray} Escursioni}
\fancyfoot[L]{\footnotesize\color{gray} Giovanni Veronese, Sebastiano Lorenzi}
\fancyfoot[R]{\footnotesize\color{gray} \thepage}
\renewcommand{\headrulewidth}{0.3pt}
\renewcommand{\footrulewidth}{0.3pt}
\let\oldheadrule\headrule
\renewcommand{\headrule}{\color{gray}\oldheadrule}
\let\oldfootrule\footrule
\renewcommand{\footrule}{\color{gray}\oldfootrule}

\hypersetup{
    colorlinks=true,
    linkcolor=black,
    urlcolor=blue,
    pdftitle={Relazione Progetto: Escursioni},
    pdfauthor={Giovanni Veronese, Sebastiano Lorenzi},
    pdfsubject={Basi di Dati},
    bookmarks=true,
    bookmarksopen=true,
    bookmarksopenlevel=2
}

\definecolor{sqlkeyword}{RGB}{0, 85, 196}
\definecolor{sqlcomment}{RGB}{0, 128, 0}
\definecolor{sqlstring}{RGB}{163, 21, 21}
\definecolor{sqlback}{RGB}{250, 251, 252}
\definecolor{sqlframe}{RGB}{224, 228, 232}
\definecolor{codegray}{rgb}{0.3,0.3,0.3}

\lstdefinestyle{sqlstyle}{
    commentstyle=\color{sqlcomment}\itshape,
    keywordstyle=\color{sqlkeyword}\bfseries,
    numberstyle=\tiny\color{codegray},
    stringstyle=\color{sqlstring},
    basicstyle=\ttfamily\footnotesize,
    breakatwhitespace=false,
    breaklines=true,
    captionpos=b,
    keepspaces=true,
    numbers=left,
    numbersep=0pt,
    showspaces=false,
    showstringspaces=false,
    showtabs=false,
    tabsize=4,
    xleftmargin=5pt
}
\lstset{style=sqlstyle, language=SQL}

\newtcblisting{sqlcode}[1][]{
    colback=sqlback,
    colframe=sqlframe,
    arc=3pt,
    boxrule=0.5pt,
    left=5pt,
    right=5pt,
    top=5pt,
    bottom=5pt,
    before skip=20pt,
    after skip=20pt,
    width=0.9\textwidth,
    center,
    listing only,
    listing options={style=sqlstyle, #1}
}

\begin{document}

% frontespizio
\begin{titlepage}
    \centering
    \vspace*{1cm}
    
    \includegraphics[width=4.5cm]{logo_unipd.png} \\ [1cm]
    
    {\scshape\LARGE Università degli Studi di Padova \par}
    \vspace{0.5cm}
    {\Large Scuola di Scienze \par}
    {\large Corso di Laurea in Informatica \par}
    
    \vspace{3.5cm}
    
    {\LARGE Progetto - Basi di Dati \par}
    \vspace{0.8cm}
    {\Huge\bfseries Escursioni \par}
    
    \vspace{3.5cm}
        \begin{minipage}[t]{0.45\textwidth}
        \begin{flushleft} \large
            \emph{Studenti:}\\
            \textbf{Giovanni Veronese}\\
            \textbf{Sebastiano Lorenzi}
        \end{flushleft}
    \end{minipage}
    \hfill
    \begin{minipage}[t]{0.45\textwidth}
        \begin{flushright} \large
            \emph{Docente:}\\
            \textbf{Prof. Massimiliano de Leoni}
        \end{flushright}
    \end{minipage}
    \vfill
    {\large \textbf{Anno Accademico:} 2025-2026 \par}
    \vspace*{1cm}
\end{titlepage}

\newpage

\tableofcontents
\thispagestyle{empty}

\newpage
\setcounter{page}{2}

\section{Abstract}
Questo progetto sviluppa una base di dati per gestire la ricerca e la partecipazione a escursioni in montagna. 
Il principale obiettivo è poter ricercare l'itinerario su misura per sé, partendo da criteri molto diversi: se si preferisce l'aiuto di un esperto si può vedere in quali date sono previste le escursioni guidate.
È pensato anche per essere un sistema che incrocia sentieri e punti d'interesse, perciò si può partire con dei luoghi in cui arrivare e trovare il percorso o al contrario dal percorso scoprire le tappe che si attraversano.
I sentieri appartengono a un parco naturale che sulla sua superficie può ospitare diversi sentieri e ancor più tappe. 
Un sentiero ha un luogo di partenza e uno di arrivo, che possono coincidere.
Uno dei punti d'interesse può essere un rifugio, affidato a un responsabile e con possibilità di avere un numero di posti letto per il pernottamento.
Le persone iscritte alla lista sono identificate dal codice fiscale, possono essere normali escursionisti o delle guide, che organizzano delle escursioni guidate in una determinata data a cui possono partecipare un certo numero di persone.

\section{Analisi dei requisiti}\label{sec:section1}
Seguono i requisiti della base di dati:

\noindent \textbf{Persona.} l'essere inserita nel sistema è già una conferma di aver superato i requisiti fisici e medici. Ogni persona è identificata da
\begin{itemize}
    \item \textbf{codice fiscale} che distingue le persone in modo univoco 
    \item \textbf{nome}
    \item \textbf{cognome} 
    \item \textbf{tipo persona} ovvero se è un escursionista o una guida
\end{itemize}

Ogni persona può partecipare a un evento organizzato come \textbf{escursionista} o come \textbf{guida}.

\noindent \textbf{Guida.} eredita tutti gli attributi di Persona, e aggiunge un attributo
\begin{itemize}
    \item \textbf{numero tesserino}, ogni guida deve essere identificata come guida abilitata tramite un codice
\end{itemize}

\noindent \textbf{Escursionista.} eredita tutti gli attributi di Persona, ma implementa l'attributo
\begin{itemize}
    \item \textbf{scadenza certificato medico} che non permette la partecipazione a una Escursione se la data di scadenza non è successiva
\end{itemize}

\noindent \textbf{Tappa.} è identificata da un codice univoco e permette di sapere i luoghi dove passa un Sentiero con i seguenti attributi:
\begin{itemize}
    \item \textbf{id tappa} per non creare ambiguità,
    \item \textbf{nome} spiega o descrive il tipo di luogo
    \item \textbf{altitudine} in metri, trattandosi di montagna è un dato utile da avere
    \item \textbf{provincia} identifica la provincia in cui si trova quel luogo
    \item \textbf{copertura di rete} un valore binario che informa la presenza (1) o meno (0) di segnale telefonico. In caso di dubbio il dato è 0.
\end{itemize}

\noindent \textbf{Rifugio.} eredita gli attributi di Tappa, e fornisce informazioni aggiuntive essendo una Tappa strategica
\begin{itemize}
    \item \textbf{id tappa} ereditata da Tappa
    \item \textbf{posti letto} indica la presenza e la possibilità di pernottamento
    \item \textbf{codice fiscale del gestore}, un rifugio inserito nell'elenco deve necessariamente avere una persona di riferimento, iscritta all'elenco Persona
    \item \textbf{proprietà} i rifugi possono essere proprietà di enti locali, privati o delle varie sezioni del Club Alpino Italiano
\end{itemize}

\noindent \textbf{Cima.} eredita gli attributi di Tappa, e fornisce una informazione in più
\begin{itemize}
    \item \textbf{id tappa} ereditata da Tappa
    \item \textbf{gruppo montuoso} indica l'appartenenza della specifica cima all'insieme di montagne dove si colloca
\end{itemize}

\noindent \textbf{Parco.} indica l'area geografica dove sarà collocato il Sentiero
\begin{itemize}
    \item \textbf{nome} è la descrizione testuale del parco Naturale o Regionale
    \item \textbf{regione} descrive la regione di appartenenza del parco
    \item \textbf{superficie} l'area espressa in $\mathrm{km}^2$ del territorio descritto
\end{itemize}

\noindent \textbf{Sentiero.} un percorso che ha un inizio e una fine che possono coincidere, descritto dai seguenti attributi
\begin{itemize}
    \item \textbf{codice sentiero} codice CAI assegnato al sentiero, univoco e utile per le indicazioni
    \item \textbf{nome} descrizione testuale del percorso
    \item \textbf{difficoltà} indica all'escursionista la difficoltà stimata del percorso
    \item \textbf{lunghezza} espressa in km, indica quanto è lungo il percorso
    \item \textbf{tempo di percorrenza} espressa in minuti, suggerisce il tempo standard di percorrenza
    \item \textbf{dislivello} espresso in metri, indica il dislivello positivo. Insieme alla lunghezza permette di farsi un'idea precisa dell'impegno necessario per una particolare escursione
    \item \textbf{numero di tappe} un valore intero che indica le tappe salvate nel database che saranno attraversate in questo percorso, senza contare il punto di partenza e il punto di arrivo
    \item \textbf{nome del parco} indica in quale parco Nazionale o Regionale appartiene il percorso. Molti sentieri non appartengono a uno specifico parco ma al generico territorio regionale
    \item \textbf{regione del parco} regione di appartenenza del percorso
    \item \textbf{id tappa di partenza} indica il luogo da cui parte il percorso
    \item \textbf{id tappa di arrivo} indica il luogo in cui il percorso finisce. Può corrispondere alla partenza in caso di percorsi ad anello
\end{itemize}

\noindent \textbf{Escursione.} è un evento a cui ci si può iscrivere e partecipare, caratterizzato dalle seguenti informazioni
\begin{itemize}
    \item \textbf{id escursione} univoco per ogni escursione programmata
    \item \textbf{data programmata} stabilisce il giorno per cui è previsto l'evento
    \item \textbf{costo} la quota necessaria per partecipare
    \item \textbf{codice sentiero} indica quale sentiero si andrà a percorrere
    \item \textbf{posti disponibili} indica la disponibilità dei posti per quello specifico evento. Se la difficoltà del sentiero è maggiore, i posti disponibili sono tendenzialmente più bassi per una maggiore sicurezza
    \item \textbf{cf guida} identifica la persona che sarà alla guida del gruppo.
\end{itemize}

\noindent La tabella \textbf{attraversa} tiene traccia di tutte le tappe collegate ai sentieri, collegando
\begin{itemize}
    \item \textbf{codice del sentiero}
    \item \textbf{id tappa}
\end{itemize}

\noindent La tabella \textbf{partecipa} crea un collegamento tra gli escursionisti e la partecipazione a ogni escursione, creando uno storico
\begin{itemize}
    \item \textbf{cf escursionista}
    \item \textbf{id escursione}
\end{itemize}

\section{Progettazione concettuale}

\begin{figure}[htbp]
    \centering
    \includegraphics[width=1\textwidth]{escursioni_finale (1).pdf}
    \caption{Diagramma E/R relativo alle escursioni.}
    \label{fig:figura1}
\end{figure}

La \autoref{fig:figura1} mostra i requisiti descritti nella \autoref{sec:section1}.

Nel database vengono rappresentate le entità principali delle escursioni: Persona, Guida, Escursionista, Tappa, Rifugio, Cima, Parco, Sentiero ed Escursione. Le relazioni tra questi concetti consentono scegliere un sentiero, anche in base alle tappe che lo compongono; le escursioni organizzate da una guida e la partecipazione degli escursionisti. La specializzazione tra Persona, Guida ed Escursionista consente di distinguere in modo chiaro i ruoli del senza duplicare informazioni. Allo stesso modo funzionano le specializzazioni di Tappa denominate Cima e Rifugio. Inoltre, la relazione tra Sentiero e Parco permette di associare ogni itinerario alla zona geografica di riferimento, mentre quella tra Rifugio e Persona definisce l'unico responsabile del luogo. Il sistema è progettato sia per gestire i dati statici sia i dati relativi all'organizzazione delle escursioni oppure fare un semplice piano di ricerca a scopo informativo.

\section{Progettazione logica}
\subsection{Analisi delle ridondanze}

L'attributo \textbf{num\_tappe} dell'entità \textbf{Sentiero}, che memorizza il numero di tappe appartenenti a un sentiero, rappresenta una ridondanza. Tale informazione può infatti essere ricavata contando il numero di tappe associate al sentiero tramite la relazione \textbf{attraversa}.

Le operazioni considerate nell'analisi sono:

\begin{itemize}
    \item \textbf{Operazione 1 (10 volte):} inserimento di una nuova tappa in un sentiero.
    \item \textbf{Operazione 2 (100 volte):} visualizzazione del numero di tappe di un sentiero.
\end{itemize}

Assumendo i seguenti volumi nella base di dati:

\begin{table}[h]
\centering
\begin{tabular}{l l r}
\toprule
Concetto & Costrutto & Volume \\
\midrule
SENTIERO & E & 100 \\
TAPPA & E & 400 \\
ATTRAVERSA & R & 1000 \\
\bottomrule
\end{tabular}
\caption{Volumi considerati per l'analisi delle ridondanze.}
\end{table}

Essendo \textbf{attraversa} una relazione molti-a-molti, è possibile che diverse tappe siano condivise da più sentieri. Per questo motivo il numero di istanze della relazione è superiore al numero delle tappe.

Si assume inoltre che il costo di una scrittura sia pari a \textbf{1}, mentre quello di una lettura sia pari a \textbf{2}.

\subsubsection{Con ridondanza}

\textbf{Operazione 1}

\begin{table}[h]
\centering
\begin{tabular}{l l c c}
\toprule
Concetto & Costrutto & Accessi & Tipo \\
\midrule
TAPPA & E & 1 & S \\
ATTRAVERSA & R & 1 & S \\
SENTIERO & E & 1 & L \\
SENTIERO & E & 1 & S \\
\bottomrule
\end{tabular}
\caption{Accessi dell'operazione 1.}
\end{table}

Operazione eseguita \textbf{10 volte}. 

\textbf{Operazione 2}

\begin{table}[h]
\centering
\begin{tabular}{l l c c}
\toprule
Concetto & Costrutto & Accessi & Tipo \\
\midrule
SENTIERO & E & 1 & L \\
\bottomrule
\end{tabular}
\caption{Accessi dell'operazione 2.}
\end{table}

Operazione eseguita \textbf{100 volte}.

\textbf{Costo totale}

\[
10 \times (3 \times 2 + 1) + 100 = 170
\]

\subsubsection{Senza ridondanza}

\textbf{Operazione 1}

\begin{table}[h]
\centering
\begin{tabular}{l l c c}
\toprule
Concetto & Costrutto & Accessi & Tipo \\
\midrule
TAPPA & E & 1 & S \\
ATTRAVERSA & R & 1 & S \\
\bottomrule
\end{tabular}
\caption{Accessi dell'operazione 1.}
\end{table}

Operazione eseguita \textbf{10 volte}.

\textbf{Operazione 2}

Poiché il numero di tappe non è memorizzato, è necessario contare le istanze della relazione \textbf{Attraversa} associate al sentiero.

Considerando:

\[
1000 / 100 = 10
\]

in media ogni sentiero attraversa \textbf{10 tappe}.

\begin{table}[h]
\centering
\begin{tabular}{l l c c}
\toprule
Concetto & Costrutto & Accessi & Tipo \\
\midrule
SENTIERO & E & 1 & L \\
ATTRAVERSA & R & 10 & L \\
\bottomrule
\end{tabular}
\caption{Accessi dell'operazione 2.}
\end{table}

Operazione eseguita \textbf{100 volte}.

\textbf{Costo totale}

\[
10 \times (2 \times 2) + 100 \times (10 + 1) = 40 + 1100 = 1140
\]


\subsubsection{Conclusioni}

L'analisi mostra che mantenere l'attributo ridondante \textbf{num\_tappe} comporta un costo maggiore nelle operazioni di aggiornamento, ma riduce notevolmente il costo delle operazioni di consultazione, che risultano molto più frequenti. Per questo motivo è conveniente mantenere l'attributo ridondante nello schema.
\subsection{Eliminazione delle generalizzazioni}

Nello schema E-R sono presenti due generalizzazioni, entrambe di tipo \textbf{parziale e disgiunta}:

\begin{itemize}
    \item \textbf{Persona}, specializzata nelle entità \textbf{Guida} ed \textbf{Escursionista};
    \item \textbf{Tappa}, specializzata nelle entità \textbf{Rifugio} e \textbf{Cima}.
\end{itemize}

Per entrambe le generalizzazioni è stata adottata la strategia di mantenere la superclasse e le relative sottoclassi.

\subsubsection{Generalizzazione Persona}

La generalizzazione dell'entità \textbf{Persona} è stata mantenuta poiché le sottoclassi \textbf{Guida} ed \textbf{Escursionista} possiedono attributi specifici (\textbf{NumeroTesserino} e \textbf{ScadenzaIscrizione}) e partecipano a relazioni differenti con l'entità \textbf{Escursione} (\textbf{Dirige} e \textbf{Partecipa}).

L'eliminazione della superclasse avrebbe comportato la duplicazione degli attributi comuni (\texttt{CF}, \textbf{Nome} e \textbf{Cognome}) nelle due sottoclassi, introducendo una ridondanza non necessaria.

\subsubsection{Generalizzazione Tappa}

Anche la generalizzazione dell'entità \textbf{Tappa} è stata mantenuta.

Le sottoclassi \textbf{Rifugio} e \textbf{Cima} possiedono attributi specifici (\textbf{posti\_letto} e \textbf{gruppo\_montuoso}) che non sono condivisi. Inoltre, essendo la generalizzazione \textbf{parziale}, una tappa può anche non appartenere ad alcuna delle due sottoclassi, rappresentando ad esempio un punto panoramico, un bivio o un'altra tappa significativa del percorso.

Il mantenimento della superclasse consente quindi di rappresentare correttamente tutte le tipologie di tappa, evitando sia la duplicazione degli attributi comuni sia la presenza di valori nulli negli attributi specifici delle sottoclassi.

In entrambe le generalizzazioni, le sottoclassi ereditano l'identificatore della rispettiva superclasse.


\subsection{Partizionamento e accorpamento di entità e relazioni}

Durante la ristrutturazione dello schema E-R è stata valutata la possibilità di accorpare alcune entità e relazioni.

Non sono stati effettuati ulteriori accorpamenti, poiché le entità presenti rappresentano concetti distinti del dominio applicativo e possiedono attributi o relazioni specifiche che ne giustificano il mantenimento.

Anche le relazioni \textbf{attraversa}, \textbf{parte\_da} e \textbf{arriva\_a} sono state mantenute separate, in quanto descrivono vincoli semantici differenti. In particolare, \textbf{parte\_da} individua la tappa iniziale di un sentiero, \textbf{arriva\_a} la tappa finale, mentre \textbf{attraversa} rappresenta l'insieme delle tappe appartenenti al percorso. Un loro accorpamento avrebbe richiesto l'introduzione di attributi discriminanti e avrebbe reso il modello meno chiaro e più complesso da gestire.

Si è pertanto scelto di mantenere inalterata la struttura delle entità e delle relazioni, preservando la chiarezza dello schema concettuale e la corretta rappresentazione del dominio.



\subsection{Scelta degli identificatori primari}

La scelta degli identificatori primari è stata effettuata privilegiando, ove possibile, attributi naturali univoci e, negli altri casi, identificatori artificiali che garantiscono l'univocità delle istanze.

\subsubsection{Persona}

L'entità \textbf{Persona} è identificata dal \textbf{Codice Fiscale (cf)}, che rappresenta un identificatore naturale univoco e stabile nel tempo.

\subsubsection{Guida ed Escursionista}

Le entità \textbf{Guida} ed \textbf{Escursionista} ereditano l'identificatore della superclasse \textbf{Persona}. Di conseguenza, il \textbf{cf} identifica univocamente anche le istanze delle due sottoclassi.

\subsubsection{Parco}

L'entità \textbf{Parco} è identificata dalla coppia \textbf{(nome, regione)}. La scelta di una chiave composta consente di distinguere univocamente i parchi, evitando possibili ambiguità nel caso di parchi con lo stesso nome situati in regioni diverse.

\subsubsection{Sentiero}

 L'entità \textbf{Sentiero} è identificata da un \textbf{codice} univoco assegnato a ciascun percorso, evitando possibili ambiguità dovute alla presenza di sentieri con lo stesso nome.

\subsubsection{Escursione}

L'entità \textbf{Escursione} è identificata da un codice \textbf{id} univoco. L'utilizzo di un identificatore evita la definizione di chiavi composte e consente di distinguere univocamente più escursioni organizzate sullo stesso sentiero.

\subsubsection{Tappa}

L'entità \textbf{Tappa} è identificata da un \textbf{ID} univoco. L'utilizzo di un identificatore artificiale evita possibili duplicazioni dovute a tappe con lo stesso nome o situate alla medesima altitudine.

\subsubsection{Rifugio e Cima}

Le entità \textbf{Rifugio} e \textbf{Cima} ereditano l'identificatore dell'entità \textbf{Tappa}, garantendo l'univocità delle rispettive istanze senza introdurre ulteriori identificatori.

\subsection{Schema Relazionale}

\begin{figure}[htbp]
    \centering
    \includegraphics[width=1\textwidth]{escursioni_ristrutturato (1).pdf}
    \caption{Diagramma E/R ristrutturato}
    \label{fig:figura2}
\end{figure}

Lo schema ristrutturato in \autoref{fig:figura2} contiene costrutti in corrispettivi dello schema relazionale. Lo schema logico è rappresentato a seguire, dove l’asterisco dopo il nome degli attributi indica quelli che ammettono valori nulli.
\begin{itemize}
    \item Persona(\underline{cf}, nome, cognome)

    \item Guida(\underline{cf}, num\_tesserino)
    \begin{itemize}
        \item Guida.Persona $\rightarrow$ Persona.cf
    \end{itemize}

    \item Escursionista(\underline{cf}, scadenza\_certificato\_medico)
    \begin{itemize}
        \item Escursionista.Persona $\rightarrow$ Persona.cf
    \end{itemize}

    \item \textbf{Parco}(\underline{nome}, \underline{regione}, superficie\_km*\textsuperscript{2})

    \item \textbf{Sentiero}(\underline{codice}, nome, difficolta, lunghezza, tempo\_percorrenza, dislivello, num\_tappe, nome\_parco, regione\_parco, partenza, arrivo)
    \begin{itemize}
        \item Sentiero.(parco\_nome, parco\_regione) $\rightarrow$ Parco.(nome, regione)
        \item Sentiero.partenza $\rightarrow$ Tappa.id
        \item Sentiero.arrivo $\rightarrow$ Tappa.id
    \end{itemize}

    \item \textbf{Tappa}(\underline{id}, nome, altitudine\_m*, provincia*, copertura\_rete\_mobile)

    \item \textbf{Rifugio}(\underline{id\_tappa}, posti\_letto, gestore, proprieta*)
    \begin{itemize}
        \item Rifugio.Tappa $\rightarrow$ Tappa.id
        \item Rifugio.gestore $\rightarrow$ Persona.cf
    \end{itemize}

    \item \textbf{Cima}(\underline{id\_tappa}, gruppo\_montuoso)
    \begin{itemize}
        \item Cima.Tappa $\rightarrow$ Tappa.ID
    \end{itemize}

    \item \textbf{Escursione}(\underline{id}, data\_programmata, costo, sentiero, guida)
    \begin{itemize}
        \item Escursione.Sentiero $\rightarrow$ Sentiero.codice
        \item Escursione.Guida $\rightarrow$ Guida.Persona
    \end{itemize}

    \item \textbf{attraversa}(\underline{Sentiero}, \underline{Tappa})
    \begin{itemize}
        \item Attraversa.Sentiero $\rightarrow$ Sentiero.codice
        \item Attraversa.Tappa $\rightarrow$ Tappa.id
    \end{itemize}

    \item \textbf{partecipa}(\underline{Escursionista}, \underline{Escursione})
    \begin{itemize}
        \item Partecipa.Escursionista $\rightarrow$ Escursionista.Persona
        \item Partecipa.Escursione $\rightarrow$ Escursione.id
    \end{itemize}
\end{itemize}

\section{Implementazione in PostgreSQL e definizione delle Query}
Il file \texttt{escursioni.sql} contiene il codice SQL necessario per la creazione e il popolamento delle tabelle del database. Questo file include inoltre una serie di Query per l’estrazione dei dati e un indice creato specificamente per migliorare le prestazioni di una di queste interrogazioni.

\subsection{Definizione delle Query}
Di seguito sono descritte le Query con i relativi output generati


\subsubsection{Query 1}
\indent Trovare il numero di escursioni organizzate da ogni guida e ordinare stampando i dati in ordine decrescente di numero di escursioni
\begin{sqlcode}[caption={Query 1}, label={}]
    SELECT g.cf, p.nome, p.cognome, COUNT(e.id_escursione) AS numero_escursioni
    FROM GUIDA g
    JOIN PERSONA p ON g.cf = p.cf
    LEFT JOIN ESCURSIONE e ON g.cf = e.cf_guida
    GROUP BY g.cf, p.nome, p.cognome
    ORDER BY numero_escursioni DESC;
\end{sqlcode}

Di seguito un estratto dell'output:

\begin{figure}[H]
    \centering
    \includegraphics[width=0.8\textwidth]{query1.png}
    \label{fig:query1}
\end{figure}


\subsubsection{Query 2}
\indent Cercare e elencare i sentieri con il massimo numero di tappe visitabili
\begin{sqlcode}[caption={Query 2}, label={}]
    DROP VIEW IF EXISTS V_NUMERO_TAPPE;

    CREATE VIEW V_NUMERO_TAPPE AS
    SELECT s.codice, s.nome, COUNT(a.id_tappa) AS numero_tappe
    FROM SENTIERO s
    JOIN attraversa a ON s.codice = a.codice_sentiero
    GROUP BY s.codice, s.nome;

    SELECT codice, nome, numero_tappe
    FROM V_NUMERO_TAPPE
    WHERE numero_tappe = (
        SELECT MAX(numero_tappe)
        FROM V_NUMERO_TAPPE
    );
\end{sqlcode}

Di seguito un l'output:

\begin{figure}[H]
    \centering
    \includegraphics[width=0.7\textwidth]{query2.png}
    \label{fig:query2}
\end{figure}


\subsubsection{Query 3}
\indent Estrarre un elenco delle escursioni comprese in un intervallo di date, con guida e sentiero, stampate in ordine cronologico
\begin{sqlcode}[caption={Query 3}, label={}]
    SELECT e.id_escursione, e.data_programmata, s.nome AS sentiero, p.nome, p.cognome, e.costo
    FROM ESCURSIONE e
    JOIN SENTIERO s ON e.codice_sentiero = s.codice
    JOIN GUIDA g ON e.cf_guida = g.cf
    JOIN PERSONA p ON g.cf = p.cf
    WHERE e.data_programmata BETWEEN '2026-06-13' AND '2026-12-24'
    ORDER BY e.data_programmata;
\end{sqlcode}

Di seguito un estratto dell'output:

\begin{figure}[H]
    \centering
    \includegraphics[width=0.9\textwidth]{query3.png}
    \label{fig:query3}
\end{figure}


\subsubsection{Query 4}
\indent Elenca gli che hanno partecipato almeno a N escursioni, ordinati per numero decrescente di escursioni effettuate
\begin{sqlcode}[caption={Query 4}, label={}]
    SELECT e.cf, p.nome, p.cognome, COUNT(pa.id_escursione) AS numero_escursioni
    FROM ESCURSIONISTA e
    JOIN PERSONA p ON e.cf = p.cf
    JOIN partecipa pa ON e.cf = pa.cf_escursionista
    GROUP BY e.cf, p.nome, p.cognome
    HAVING COUNT(pa.id_escursione) >= 6
    ORDER BY numero_escursioni DESC;
\end{sqlcode}

Di seguito un estratto dell'output:

\begin{figure}[H]
    \centering
    \includegraphics[width=0.7\textwidth]{query4.png}
    \label{fig:query4}
\end{figure}


\subsubsection{Query 5}
\indent Cerca e calcola il costo medio delle escursioni organizzate da ogni guida, ordinati per costo decrescente

\begin{sqlcode}[caption={Query 5}, label={}]
    SELECT g.cf, p.nome, p.cognome, ROUND(AVG(e.costo),2) AS costo_medio
    FROM GUIDA g
    JOIN PERSONA p ON g.cf = p.cf
    JOIN ESCURSIONE e ON g.cf = e.cf_guida
    GROUP BY g.cf, p.nome, p.cognome
    ORDER BY costo_medio DESC;
\end{sqlcode}

Di seguito un estratto dell'output:

\begin{figure}[H]
    \centering
    \includegraphics[width=0.7\textwidth]{query5.png}
    \label{fig:query5}
\end{figure}


\subsection{Creazione degli indici}
Per ottimizzare la Query è stato scelto un indice di tipo B-tree sull'attributo \textbf{data\_programmata} della tabella \texttt{ESCURSIONE}. Questa scelta è motivata dal fatto che la Query effettua una ricerca su un intervallo di date mediante l'operatore \texttt{BETWEEN} e restituisce i risultati ordinati cronologicamente tramite \texttt{ORDER BY}. La struttura B-tree mantiene i valori dell'indice in ordine, consentendo di individuare rapidamente l'inizio dell'intervallo richiesto e di scorrere solo le tuple comprese tra le due date, evitando una scansione completa della tabella. Poiché i dati sono già ordinati nell'indice, anche l'operazione di ordinamento risulta più efficiente. Un indice di tipo \textbf{Hash} non sarebbe invece adatto, poiché è progettato esclusivamente per ricerche basate sull'uguaglianza (\texttt{=}) e non supporta in modo efficiente ricerche per intervallo o operazioni di ordinamento. Per questi motivi, il B-tree rappresenta la scelta più appropriata per questa query. L'indice descritto nelle righe precedenti è questo:

\begin{sqlcode}[caption={Query 5}, label={}]
    CREATE INDEX idx_escursione_data ON ESCURSIONE(data_programmata);
\end{sqlcode}

\section{Applicazione software}

Il file \texttt{codice.c} contiene il codice \texttt{C} necessario per connettersi al database e visualizzare a schermo i risultati delle Query presentate nella Sezione 5. Una volta compilato ed eseguito il programma, verrà mostrato all’utente un menu con la seguente lista di interrogazioni disponibili:
\begin{enumerate}
    \item Numero di escursioni organizzate da ogni guida
    \item Sentieri con il massimo numero di tappe visitabili
    \item Elenco delle escursioni comprese in un intervallo di date, con guida e sentiero, ordinate cronologicamente
    \item Escursionisti che ahanno partecipato almeno a N escursioni
    \item Costo medio delle escursioni organizzate da ogni guida
\end{enumerate}

L’utente potrà selezionare la Query da eseguire digitando il numero corrispondente. In particolare, la Query 4, al momento della selezione, il programma richiederà l’inserimento di un numero N che è il numero minimo di escursioni da adottare come criterio di ricerca.
\end{document}